%%%%%%%%%%%%%%%%%%%%%%%%%%%%%%%%%%%%%%%%%%%%%%%%%%%%%%%%%%%%%%%%%%%%%%%%%%%%%%
% Fiche d'évaluation AMC — Projet cellule robotisée (TPs 2 à 5)
%
% Évalue la configuration complète de l'automate M340 et la programmation
% UniVALplc sous EtherNet/IP : configuration NOC, structure logicielle,
% GEMMA, modes de marche, qualité du code et investissement.
%
% Barème : 20 pts
%   - Configuration EtherNet/IP   : 3 pts  (1 critère)
%   - Structure logicielle        : 4 pts  (2 critères)
%   - GEMMA et modes de marche    : 6 pts  (3 critères)
%   - Qualité globale du code     : 2 pts  (1 critère)
%   - Investissement dans les séances : 5 pts (2 critères)
%
% Compilation : pdflatex fiche_projet.tex
%%%%%%%%%%%%%%%%%%%%%%%%%%%%%%%%%%%%%%%%%%%%%%%%%%%%%%%%%%%%%%%%%%%%%%%%%%%%%%

\documentclass[QCM, noCustomPackages]{UPSTI_Document}

\UPSTIloadProfile{BUT_GEII}

% ============================================================
% MÉTADONNÉES
% ============================================================
\discipline[Autom - Robotique]{Automatisme pour la robotique}
\newcommand{\UPSTItitreEnTete}{Évaluation projet}
\newcommand{\UPSTItitre}{Cellule robotisée UniVALplc --- Évaluation projet}
\newcommand{\UPSTInumero}{Projet}
\newcommand{\sequence}{99}
\documentVersion{E}
\newcommand{\UPSTInumeroVersion}{--}
\newcommand{\UPSTImessage}{}

% ============================================================
\begin{document}

\AMCrandomseed{202500}
\setdefaultgroupmode{withoutreplacement}

% ============================================================
% BANQUE DE CRITÈRES
% ============================================================

% ------------------------------------------------------------
% Configuration EtherNet/IP — NOC_ETHIP_ROBOT (3 pts)
% ------------------------------------------------------------

\element{config}{%
  \UPSTIcritere[3]{noc-robot}{%
    Configuration complète du réseau \texttt{NOC\_ETHIP\_ROBOT} :
    équipements CS9 et M221 ajoutés dans le bon ordre, adresses IP correctes,
    tailles des données E/S configurées (Exclusive Owner / assemblage).%
  }%
}

% ------------------------------------------------------------
% Structure logicielle de base (4 pts)
% ------------------------------------------------------------

\element{structure}{%
  \UPSTIcritere[2]{struct-base}{%
    Objet \texttt{T\_StaeubliRobot} réservé ; blocs \texttt{VAL\_ReadAxesGroup}
    et \texttt{VAL\_WriteAxesGroup} appelés aux bons moments du cycle API
    (lecture en début, écriture en fin).%
  }%
}

\element{structure}{%
  \UPSTIcritere[2]{vars-echange}{%
    Variables d'échange robot et HMI déclarées avec les bons types
    et correctement câblées aux blocs de lecture/écriture.%
  }%
}

% ------------------------------------------------------------
% GEMMA et modes de marche (6 pts)
% ------------------------------------------------------------

\element{gemma}{%
  \UPSTIcritere[2]{gemma-securite}{%
    Modes \textbf{D1} et \textbf{D2} fonctionnels : arrêt d'urgence
    (bras hors puissance) et reprise après défaillance
    (reset, vidage pile, remise sous puissance).%
  }%
}

\element{gemma}{%
  \UPSTIcritere[2]{gemma-init}{%
    Mode \textbf{A6} fonctionnel : mise en conditions initiales du robot
    (séquence de déplacement conforme, position finale correcte).%
  }%
}

\element{gemma}{%
  \UPSTIcritere[2]{mode-f1}{%
    Mode \textbf{F1} Production Normale fonctionnel : cycle robot + guichet
    coordonnés, macro-étapes correctement définies.%
  }%
}

% ------------------------------------------------------------
% Qualité globale du code (2 pts)
% ------------------------------------------------------------

\element{qualite}{%
  \UPSTIcritere[2]{qualite-glob}{%
    Qualité et cohérence globale du projet : nommage conforme aux
    conventions, structure lisible, absence de code mort.%
  }%
}

% ------------------------------------------------------------
% Investissement dans les séances (5 pts)
% ------------------------------------------------------------

\element{investissement}{%
  \UPSTIcritere[2]{invt-presence}{%
    Présence et engagement : travail effectif durant les séances,
    sollicitation appropriée de l'enseignant.%
  }%
}

\element{investissement}{%
  \UPSTIcritere[3]{invt-autonomie}{%
    Progression et autonomie : évolution observée sur la durée du projet,
    capacité à analyser et résoudre les problèmes sans aide systématique.%
  }%
}

% ============================================================
% GÉNÉRATION DES COPIES NOMINATIVES
% ============================================================

\csvreader[head to column names]{etudiants.csv}{}{%
  \exemplaire{1}{%

    \UPSTIresetCritere
    \renewcommand{\UPSTImessage}{\champnom{\nom\ \prenom}}%
    \renewcommand{\UPSTInumeroVersion}{\id{}}%
    \AMCassociation{\id}
    \UPSTIbuildPage

    % --------------------------------------------------------
    % Configuration EtherNet/IP (3 pts)
    % --------------------------------------------------------
    \section*{Configuration EtherNet/IP
              \hfill \normalfont\small /3~pts}

    \restituegroupe{config}

    \smallskip
    \noindent\textit{Commentaires :}\\[2pt]
    \UPSTIquadrillage{2}

    \bigskip

    % --------------------------------------------------------
    % Structure logicielle de base (4 pts)
    % --------------------------------------------------------
    \section*{Structure logicielle de base
              \hfill \normalfont\small /4~pts}

    \restituegroupe{structure}

    \smallskip
    \noindent\textit{Commentaires :}\\[2pt]
    \UPSTIquadrillage{2}

    \bigskip

    % --------------------------------------------------------
    % GEMMA et modes de marche (6 pts)
    % --------------------------------------------------------
    \section*{GEMMA et modes de marche
              \hfill \normalfont\small /6~pts}

    \restituegroupe{gemma}

    \smallskip
    \noindent\textit{Commentaires :}\\[2pt]
    \UPSTIquadrillage{3}

    \bigskip

    % --------------------------------------------------------
    % Qualité globale du code (2 pts)
    % --------------------------------------------------------
    \section*{Qualité globale du code
              \hfill \normalfont\small /2~pts}

    \restituegroupe{qualite}

    \smallskip
    \noindent\textit{Commentaires :}\\[2pt]
    \UPSTIquadrillage{2}

    \bigskip

    % --------------------------------------------------------
    % Investissement dans les séances (5 pts)
    % --------------------------------------------------------
    \section*{Investissement dans les séances
              \hfill \normalfont\small /5~pts}

    \restituegroupe{investissement}

    \smallskip
    \noindent\textit{Observations :}\\[2pt]
    \UPSTIquadrillage{4}

    \AMCcleardoublepage
  }%
}

\end{document}
